\documentclass[conference]{IEEEtran}
\IEEEoverridecommandlockouts
\usepackage{cite}
\usepackage{amsmath,amssymb}
\usepackage{graphicx}
\usepackage{booktabs}
\usepackage{textcomp}
\usepackage{xcolor}
\usepackage[hidelinks]{hyperref}
\hypersetup{pdfauthor={},pdfsubject={},pdfkeywords={},pdfcreator={}}
\graphicspath{{../results/}{./}}
% Generated by `python -m flypath report` from results/*.json. Do not edit.
\newcommand{\NullB}{200}
\newcommand{\BootR}{200}
\newcommand{\NItems}{4{,}000}
\newcommand{\PrimaryK}{92}
\newcommand{\NGlomFull}{51}
\newcommand{\NCellsFull}{1{,}886}
\newcommand{\NGlom}{35}
\newcommand{\NCells}{1{,}838}
\newcommand{\NNZ}{7{,}110}
\newcommand{\InputsFull}{5.28}
\newcommand{\Inputs}{3.87}
\newcommand{\PNPartners}{5.59}
\newcommand{\FanMin}{19}
\newcommand{\FanMax}{473}
\newcommand{\FanMaxGlom}{DC1}
\newcommand{\NOdorants}{172}
\newcommand{\MissingPct}{28.7\%}
\newcommand{\DoorCommit}{db323a4}
\newcommand{\PrimaryReal}{0.4063}
\newcommand{\PrimaryNullMean}{0.4148}
\newcommand{\PrimaryNullSD}{0.0035}
\newcommand{\PrimaryNullCV}{0.85\%}
\newcommand{\PrimaryNullSE}{0.00025}
\newcommand{\PrimaryRel}{-2.06\%}
\newcommand{\PrimaryZ}{-2.42}
\newcommand{\PrimaryPtwo}{0.020}
\newcommand{\PrimaryPlower}{0.010}
\newcommand{\PrimaryRank}{200}
\newcommand{\BootLo}{-3.1\%}
\newcommand{\BootHi}{-0.8\%}
\newcommand{\EquivPrimary}{holds}
\newcommand{\MarginPrimary}{5\%}
\newcommand{\MdeUp}{2.5\%}
\newcommand{\MdeDown}{2.5\%}
\newcommand{\MinHolm}{1.000}
\newcommand{\SecPMin}{0.338}
\newcommand{\SecPMax}{0.468}
\newcommand{\SecRelMin}{-3.0\%}
\newcommand{\SecRelMax}{-0.9\%}
\newcommand{\EquivTwoFive}{does not hold}
\newcommand{\NullsAbove}{199}
\newcommand{\TieConflictMax}{2.3\%}
\newcommand{\QRealP}{0.005}
\newcommand{\QRealRel}{+3.7\%}
\newcommand{\PlateauSweeps}{1}
\newcommand{\JaccardPlateau}{10}
\newcommand{\JaccardFinal}{0.088}
\newcommand{\CovNew}{100\%}
\newcommand{\CovNewLo}{96}
\newcommand{\CovNewHi}{100}
\newcommand{\CovOld}{87\%}
\newcommand{\CovOldLo}{79}
\newcommand{\CovOldHi}{92}
\newcommand{\CovDatasets}{100}
\newcommand{\CovItems}{600}
\newcommand{\RobNNeg}{24}
\newcommand{\RobNAtMin}{13}
\newcommand{\RobPFloor}{0.049}
\newcommand{\ImputeRelMin}{-0.7\%}
\newcommand{\ImputeRelMax}{+0.6\%}
\newcommand{\ImputePMin}{0.439}
\newcommand{\RobN}{26}
\newcommand{\RobB}{40}
\newcommand{\RobRelMin}{-3.9\%}
\newcommand{\RobRelMax}{+0.8\%}
\newcommand{\RobNSig}{13}
\newcommand{\RobMinP}{0.049}
\newcommand{\ArchKbitsRatio}{1.37}
\newcommand{\ArchStorageRatio}{0.90}
\newcommand{\ArchComputeRatio}{1.33}
\newcommand{\StorageBitsPrimary}{523}
\newcommand{\ComputeBits}{102}
\newcommand{\WiringRows}{8 & 0.1973 & 0.2034 $\pm$ 0.0085 & -3.0 & 0.468 & 1.000 & +3.7 & +26.7 & +3.2 \\
16 & 0.2879 & 0.2933 $\pm$ 0.0062 & -1.8 & 0.438 & 1.000 & +2.7 & +9.7 & +3.2 \\
32 & 0.3531 & 0.3575 $\pm$ 0.0045 & -1.2 & 0.338 & 1.000 & +3.3 & +3.0 & +3.6 \\
64 & 0.3963 & 0.4001 $\pm$ 0.0039 & -0.9 & 0.348 & 1.000 & +3.4 & +0.7 & +3.4 \\
92$^\ast$ & 0.4063 & 0.4148 $\pm$ 0.0035 & -2.1 & 0.020 & -- & +3.2 & +0.1 & +3.2 \\}
\newcommand{\ArchRows}{8 & 0.197 & 0.019 & 0.264 (72) & 0.305 (102) \\
16 & 0.288 & 0.069 & 0.340 (130) & 0.305 (102) \\
32 & 0.353 & 0.151 & 0.394 (229) & 0.305 (102) \\
64 & 0.396 & 0.249 & 0.436 (397) & 0.305 (102) \\
92 & 0.406 & 0.296 & 0.454 (523) & 0.305 (102) \\}
\newcommand{\RobRows}{Baseline & -2.2 & 0.049 \\
Left hemisphere & -1.1 & 0.244 \\
$\geq$5-synapse threshold & -2.3 & 0.098 \\
Missing data (3 imputations) & -0.7 to +0.6 & 0.439--0.780 \\
Coverage thresholds (2) & -1.7 to -1.4 & 0.049--0.341 \\
All 46 glomeruli with data, imputed & -0.0 & 1.000 \\
Mixture seeds (4) & -2.4 to -1.6 & 0.049--0.244 \\
Mixture size 2 / 2--10 & -2.0 to -0.9 & 0.146--0.195 \\
Measured odorants only & -3.9 & 0.049 \\
Normalised / angular metric & -2.4 to -1.7 & 0.049 \\
Tie rules (index, include, 2 seeds) & -2.4 to -2.1 & 0.049 \\
Synapse-weighted & -1.1 & 0.488 \\
Synthetic inputs (3) & -3.1 to +0.8 & 0.049--0.244 \\}
\newcommand{\ReplRows}{SIFT & 4 & 0.174 & 0.180 & 0.283 & 0.357 & 2.1$\times$ \\
SIFT & 16 & 0.499 & 0.500 & 0.576 & 0.608 & 1.2$\times$ \\
GloVe & 4 & 0.055 & 0.056 & 0.081 & 0.093 & 1.7$\times$ \\
GloVe & 16 & 0.134 & 0.135 & 0.150 & 0.153 & 1.1$\times$ \\
MNIST & 4 & 0.159 & 0.165 & 0.272 & 0.365 & 2.3$\times$ \\
MNIST & 16 & 0.516 & 0.516 & 0.583 & 0.622 & 1.2$\times$ \\}
\newcommand{\ReplMnistLsh}{0.159}
\newcommand{\ReplMnistFly}{0.365}
\newcommand{\ReplMnistRatio}{2.3}
\newcommand{\ReplOpsBitsMnist}{390}
\newcommand{\ReplSiftRandom}{0.180}
\newcommand{\ReplSiftWta}{0.283}
\newcommand{\ReplTrials}{50}
\newcommand{\ReplOpsMin}{0.67}
\newcommand{\ReplOpsMax}{0.97}
\newcommand{\ReplFlyMax}{0.73}
\newcommand{\ConnRows}{SIFT & 4 & 0.374 & 0.375 & -0.1 & 0.784 & 0.407 & 0.197 \\
SIFT & 16 & 0.567 & 0.567 & +0.1 & 0.824 & 0.572 & 0.573 \\
SIFT & 94 & 0.722 & 0.728 & -0.8 & 0.039 & 0.776 & 0.902 \\
GloVe & 4 & 0.179 & 0.178 & +0.3 & 0.471 & 0.207 & 0.094 \\
GloVe & 16 & 0.312 & 0.313 & -0.3 & 0.353 & 0.318 & 0.316 \\
GloVe & 94 & 0.466 & 0.475 & -1.9 & 0.039 & 0.510 & 0.777 \\
MNIST & 4 & 0.445 & 0.446 & -0.2 & 0.471 & 0.470 & 0.215 \\
MNIST & 16 & 0.672 & 0.675 & -0.4 & 0.039 & 0.681 & 0.654 \\
MNIST & 94 & 0.805 & 0.810 & -0.7 & 0.039 & 0.834 & 0.935 \\}
\newcommand{\ConnB}{50}
\newcommand{\ConnTrials}{20}
\newcommand{\ConnGlom}{51}
\newcommand{\ConnCells}{1{,}886}
\newcommand{\ConnBigK}{94}
\newcommand{\ConnRelMin}{-1.9\%}
\newcommand{\ConnRelMax}{+0.3\%}
\newcommand{\ConnPMin}{0.039}
\newcommand{\ConnNSig}{6}
\newcommand{\ConnNTests}{18}
\newcommand{\ConnOpsBits}{98}
\newcommand{\ConnOpsMin}{0.78}
\newcommand{\ConnOpsMax}{0.94}
\newcommand{\ConnFlyBigMin}{0.47}
\newcommand{\ConnFlyBigMax}{0.80}
\newcommand{\ConnRandGainMax}{20\%}
\newcommand{\ConnMnistReal}{0.445}
\newcommand{\ConnMnistLsh}{0.215}
\newcommand{\ConnMnistRatio}{2.1}


\begin{document}

\title{A Connectome Test of the Fly Hashing Algorithm}

\author{\IEEEauthorblockN{Sebastian Senge}
\IEEEauthorblockA{\textit{Independent Researcher}}}

\maketitle

\begin{abstract}
The \textit{Drosophila} olfactory circuit has been modelled as a
locality-sensitive hash: glomerular input is expanded through a sparse binary
projection onto Kenyon cells and sparsified by winner-take-all. The projection
in that model was drawn at random. We replace it with the measured
glomerulus-to-Kenyon-cell connectivity of the MaleCNS connectome, released in
2026, and evaluate nearest-neighbour retrieval on synthetic linear mixtures of
partially observed DoOR odorant profiles. The measured wiring is structured:
glomerulus co-occurrence exceeds that of all \NullB{} uniform
degree-preserving rewirings. Its retrieval is slightly worse than theirs: at
the pre-specified hash size it scores below \NullsAbove{} of the \NullB{}
($\PrimaryRel$, two-sided $p=\PrimaryPtwo$; $90\%$ bootstrap interval
\BootLo{} to \BootHi{}), inside a pre-specified $\pm\MarginPrimary$
equivalence margin. The deficit shrinks to between \ImputeRelMin{} and
\ImputeRelMax{} when unmeasured responses are imputed rather than set to zero,
so it cannot be separated from the missing-data model. Against Gaussian sign
codes the fly hash leads only when codes are compared at $k$ bits; it trails
at matched storage, and at matched computation it trails for small $k$ and is
near parity for large $k$.
\end{abstract}

\begin{IEEEkeywords}
connectomics, locality-sensitive hashing, similarity search, null models
\end{IEEEkeywords}

\section{Introduction}

Locality-sensitive hashing (LSH) maps items to short codes under which similar
items collide with high probability \cite{indyk1998,charikar2002}. Dasgupta,
Stevens and Navlakha \cite{dasgupta2017} observed that the fly olfactory
circuit computes a related transform that departs from classical practice: a
sparse binary projection that expands rather than reduces dimensionality,
followed by winner-take-all. At matched computational cost it outperformed
classical LSH on standard benchmarks. The projection was sampled at random,
anchored to biology only by the number of inputs per Kenyon cell, following the
then-prevailing account that this connectivity is effectively random
\cite{caron2013}.

That account has since been revised. Zheng \textit{et al.} \cite{zheng2022}
mapped projection-neuron input to Kenyon cells in a whole-brain electron
microscopy volume, found it structured, and reported that simulations of the
observed network discriminated odours worse than a random network, except when
signal was restricted to overconvergent, largely food-responsive inputs. The
MaleCNS connectome \cite{malecns2026} now provides a complete central nervous
system of a second animal. We ask a narrower question than
\cite{zheng2022}: on the retrieval benchmark that motivated the hashing
interpretation, does substituting measured connectivity for the random
projection change retrieval quality? Our contributions are:

\begin{itemize}
\item the fly hashing benchmark evaluated under MaleCNS connectivity against
a uniform fixed-margin null, which isolates the pairing of glomeruli with
Kenyon cells from both degree sequences;
\item a comparison with Gaussian sign codes at three explicit budgets (code
length, storage, computation);
\item an inference protocol stated in advance of this analysis: an exact
randomization test with a pre-specified primary hash size, a two-stage
bootstrap whose coverage we check by simulation, a pre-specified equivalence
margin, and a power analysis.
\end{itemize}

\section{Methods}

\subsection{Projection}

Let $M\in\{0,1\}^{d\times m}$ with $M_{ji}=1$ iff glomerulus $j$ connects to
Kenyon cell $i$. We build $M$ from MaleCNS v1.0 using every reconstructed
connection, with no synapse-count threshold. Projection neurons of class
\texttt{ALPN} are assigned to glomeruli by cell type; multiglomerular,
unnamed, thermo- and hygrosensory types are excluded. For the right
hemisphere this gives $d=\NGlomFull$ glomeruli and $m=\NCellsFull$ Kenyon cells.
The odour data cover $\NGlom$ of these glomeruli, and all analyses use that
restriction: $d=\NGlom$, $m=\NCells$, $\NNZ$ non-zero entries.

We report \emph{distinct glomerular inputs} per cell, $c_i=\sum_j M_{ji}$
(mean $\Inputs$; $\InputsFull$ before restriction), and \emph{fan-out}
$f_j=\sum_i M_{ji}$ (range $\FanMin$--$\FanMax$, maximal for \FanMaxGlom). Because
projection neurons are aggregated by glomerulus, $c_i$ is not an anatomical
claw count; the number of distinct projection-neuron partners is larger (mean
$\PNPartners$), and neither quantity is validated against claw morphology here.

\subsection{Hash, baselines and budgets}

As in \cite{dasgupta2017}, an input $x\in\mathbb{R}^d_{\ge0}$ is normalised to
unit mean, $\tilde x = x/\bar x$, expanded as $y=M^\top\tilde x$, and the
$k$ largest entries of $y$ are set to one. Cells with identical inputs tie
exactly. Our rule selects exactly $k$ winners, breaking ties at the threshold
by a fixed random priority over cells; the lowest-index rule and a rule
admitting every tied cell are sensitivity analyses. Boundary ties affect at
most \TieConflictMax{} of queries at any $k$.

The baseline is a Gaussian sign code, $\mathrm{sign}(\tilde x^\top W)$ with
$W\in\mathbb{R}^{d\times b}$ i.i.d.\ normal. A $k$-sparse tag over $m$ cells
is not a $k$-bit code, so we compare at three budgets: $b=k$ bits (code length
equal to the number of active cells); $b=\lceil\log_2\binom{m}{k}\rceil$, the
length of an optimal code for an arbitrary $k$-of-$m$ tag (storage-matched,
\StorageBitsPrimary{} bits at the primary $k$); and $b=\mathrm{nnz}(M)/d=\ComputeBits$,
equal numbers of input-weight operations per item (computation-matched;
normalisation, selection and thresholding excluded on both sides).

\subsection{Null and comparison matrices}

Only $M$ varies between conditions. The null of record is uniform over binary
matrices with the same row and column sums as $M$, sampled by curveball trades
\cite{strona2014}, which Carstens proved converge to the uniform distribution
\cite{carstens2015}. We use $30$ trade attempts per cell. In five chains
run to $100$ sweeps, the glomerulus co-occurrence statistic
$Q=\sum_{j<l}(MM^\top)_{jl}^2$ stays within three null SDs from
\PlateauSweeps{} sweep on, the Jaccard overlap with the measured matrix levels
off at \JaccardFinal{} by \JaccardPlateau{} sweeps, and retrieval scores at $30$
and $100$ sweeps are indistinguishable (supplement).
Comparison matrices are not nulls and each changes something the null keeps:
input counts kept but glomeruli uniform (fan-out flattened); six uniform
inputs per cell (the construction of \cite{dasgupta2017}); input counts kept
with fan-out made uniform and pairing random; and a degraded matrix in which
all cells draw from four glomeruli.

\subsection{Benchmark}

Responses are the DoOR 2.0 consensus matrix \cite{munch2016} (commit
\texttt{\DoorCommit}), mapped from receptor to glomerulus. Keeping glomeruli
measured for $\ge 40$ odorants and odorants measured on $\ge 12$ of them leaves
$\NOdorants$ odorants, of which \MissingPct{} of entries are unmeasured. These
are set to zero in the baseline; unmeasured is not the same as non-responsive,
so three imputations are sensitivity analyses. The benchmark consists of
$\NItems$ \emph{synthetic linear mixtures of partially observed DoOR profiles}
(Dirichlet weights over two to five sources). They are not measured mixture
responses and add no biological measurements; they characterise this
generator. Ground truth is the $\kappa=10$ nearest items by Euclidean distance
on raw responses, a target the concentration-invariant hash cannot fully
recover; distances on normalised input and angular distance are sensitivity
analyses. Items are ranked by Hamming distance with ties broken by a fixed
random order per query, identical across matrices, and scored by mean average
precision (mAP).

\subsection{Inference}

The primary hash size, fixed before this analysis, is $k=\PrimaryK$, i.e.\
$5\%$ of Kenyon cells as in \cite{dasgupta2017}; $k\in\{8,16,32,64\}$ are
secondary and Holm-adjusted. Two questions are separated. \emph{Is the
measured matrix unusual among uniform fixed-margin matrices?} Under that null
hypothesis its score is exchangeable with $B=\NullB$ null scores, so the rank
gives an exact randomization $p$-value conditional on the benchmark.
\emph{How large is the difference, across odorant sets?} The estimand is
$\theta=\mathbb{E}_D[s_D(M)-\mathbb{E}_0 s_D]/\mathbb{E}_D\mathbb{E}_0 s_D$,
where $D$ is a benchmark generated from the source odorants and $\mathbb{E}_0$
is over null matrices. A two-stage bootstrap resamples source odorants,
regenerates mixtures and ground truth, recomputes retrieval, and re-estimates
the null mean from $20$ freshly drawn null matrices ($\BootR$ replicates), so
the interval carries benchmark and finite-null uncertainty. We check its
coverage by simulation from a known population fitted to the data. We
distinguish variation across random networks (null SD) from uncertainty in the
null mean (SD$/\sqrt{B}$).

Equivalence is assessed at a margin of $\pm\MarginPrimary$ relative mAP, fixed
before this analysis but after an earlier version had been seen; results at
$\pm2.5\%$ and $\pm10\%$ are in the supplement. The margin is anchored to the
effects at issue in \cite{dasgupta2017}, improvements of up to $3\times$: a
wiring effect below it could not overturn the modelling assumption being
tested. Power is estimated under a location-shift model using the empirical
null distribution.

\section{Results}

\begin{figure*}[t]
\centerline{\includegraphics[width=0.82\textwidth]{flyhash.png}}
\caption{\textbf{Left:} mAP of the fly hash on measured wiring and of Gaussian
sign codes at three budgets (mean of 20 seeds). \textbf{Right:} each
orange point is one of \NullB{} uniform fixed-margin rewirings, blue diamonds
the measured wiring, both relative to the null mean; two-sided randomization
$p$ shown.}
\label{fig:main}
\end{figure*}

\begin{table}[t]
\caption{Measured wiring against the null ($B=\NullB$) and against
comparison matrices (relative to the null mean, \%). $^\ast$primary.}
\label{tab:wiring}
\centering
\scriptsize
\setlength{\tabcolsep}{2.6pt}
\begin{tabular}{rcccccccc}
\toprule
$k$ & real & null & $\Delta$\% & $p$ & $p_\text{Holm}$ & fan flat & 6 unif. & balanced \\
\midrule
\WiringRows
\bottomrule
\end{tabular}
\end{table}

\subsection{The measured wiring is structured}

Glomerulus co-occurrence $Q$ of the measured matrix is \QRealRel{} above the
null mean and exceeds all \NullB{} null draws (upper-tail $p=\QRealP$). The
pairing of glomeruli within Kenyon cells is therefore not uniform given the
degrees, in agreement with \cite{zheng2022} in a second animal.

\subsection{Retrieval is slightly below the null}

At the primary size the measured matrix scores $\PrimaryReal$ against a null
of $\PrimaryNullMean\pm\PrimaryNullSD$ (SD across networks; SE of the mean
$\PrimaryNullSE$), a relative difference of \PrimaryRel{}. It lies below
\NullsAbove{} of the \NullB{} null draws: two-sided randomization
$p=\PrimaryPtwo$. At the secondary sizes the differences are \SecRelMin{} to
\SecRelMax{}, with unadjusted $p$ from \SecPMin{} to \SecPMax{} and none
significant after Holm adjustment (Table~\ref{tab:wiring},
Fig.~\ref{fig:main}). The sign is the same at every size, but the sizes share
one benchmark and one matrix, so this is not independent evidence.

The two-stage $90\%$ interval for $\theta$ at the primary size is \BootLo{} to
\BootHi{}: the deficit is unlikely to be zero, and equally unlikely to exceed a
few percent. Equivalence at the pre-specified $\pm\MarginPrimary$ margin
\EquivPrimary; at $\pm2.5\%$ it \EquivTwoFive. In simulation the procedure
covered the true value in \CovNew{} of \CovDatasets{} datasets (Wilson $95\%$:
\CovNewLo--\CovNewHi\%) at nominal $90\%$, so it is, if anything,
conservative; the group bootstrap of an earlier version of this analysis
covered it in \CovOld{} (\CovOldLo--\CovOldHi\%). The randomization test has
$80\%$ power against shifts of \MdeDown{} downward and \MdeUp{} upward.

\subsection{The architecture comparison depends on the budget}

\begin{table}[t]
\caption{mAP of the fly hash and of Gaussian sign codes (bits in parentheses).}
\label{tab:arch}
\centering
\scriptsize
\setlength{\tabcolsep}{3pt}
\begin{tabular}{rcccc}
\toprule
$k$ & fly & Gauss.\ $k$ bits & storage-matched & computation-matched \\
\midrule
\ArchRows
\bottomrule
\end{tabular}
\end{table}

At $b=k$ bits the fly hash is ahead at every size (Table~\ref{tab:arch}), by
a factor that falls with $k$ ($\ArchKbitsRatio\times$ at the primary size). At
matched storage it is behind at every size ($\ArchStorageRatio\times$). At
matched computation it is behind for small $k$ and slightly ahead from
$k=64$ ($\ArchComputeRatio\times$ at the primary size). The same ordering holds
under normalised and angular ground truth. The $3\times$ advantage reported in
\cite{dasgupta2017} was at matched computation on other benchmarks and larger
inputs; we do not reproduce its magnitude at the circuit's own scale.

\subsection{Robustness}

\begin{table}[t]
\caption{Primary size under one-factor variations ($B=\RobB$ nulls each):
relative difference to the null mean (\%) and two-sided $p$.}
\label{tab:rob}
\centering
\scriptsize
\setlength{\tabcolsep}{3pt}
\begin{tabular}{lcc}
\toprule
variation & $\Delta$\% & $p$ \\
\midrule
\RobRows
\bottomrule
\end{tabular}
\end{table}

Across \RobN{} variations (Table~\ref{tab:rob}) the relative difference at the
primary size ranges from \RobRelMin{} to \RobRelMax{} and is negative in
\RobNNeg{}. In \RobNAtMin{} the measured matrix falls below all \RobB{} nulls,
the smallest attainable $p$ ($\RobPFloor$); the variations are not independent,
and no adjustment is applied. The exception that matters is missing data:
imputing unmeasured responses from glomerulus means, odorant means or a
rank-5 fit instead of zero gives \ImputeRelMin{} to \ImputeRelMax{}, with
$p\geq\ImputePMin$. Every variation, with all secondary sizes, is in the
supplement.

\section{Discussion}

The measured pairing is structured, and under the baseline treatment of the
data it retrieves slightly worse than uniform rewirings with the same degrees.
The direction agrees with \cite{zheng2022}, whose observed network
discriminated odours worse than a random one. The size of the effect, about
two percent with an interval reaching three, is small beside the architectural
effects at issue in \cite{dasgupta2017}, and the pre-specified equivalence test
holds at $\pm\MarginPrimary$. For this task the random-matrix idealisation
therefore loses nothing that the measured pairing provides, and may gain a
little.

Two cautions limit that reading. First, the deficit depends on treating
unmeasured DoOR entries as zero; with any of three imputations it is no longer
detectable. The wiring and the missing-data model interact, and this
experiment cannot say which imputation is closer to the truth. Second, the
comparison matrices answer different questions from the null because each
changes a degree sequence. Flattening fan-out while keeping input counts
raises mAP above the null at every size (Table~\ref{tab:wiring}); that is an
intervention on fan-out, and it does not explain the measured matrix's
position relative to its own null, which holds fan-out fixed. That the
structure serves other ends, such as innate valence or the prioritised
channels of \cite{zheng2022}, is a hypothesis this experiment does not test.

\section{Limitations}

One animal, one hemisphere as the primary analysis, one task and one metric.
The benchmark is generated: mixtures are linear and synthetic, \MissingPct{}
of source entries are unmeasured, and $35$ of $51$ glomeruli are covered.
The model is binary and rate-free, and ignores inhibition beyond
winner-take-all. Input-dependence of the kind reported in \cite{zheng2022} is
untested, as is novelty detection \cite{dasgupta2018}.

\section{Conclusion}

In a binary glomerulus-to-Kenyon-cell model evaluated on synthetic mixtures
derived from DoOR responses, measured connectivity produced slightly lower
retrieval scores than the mean of uniform degree-preserving randomizations,
although its pairing is itself non-random. The deficit is small, within a
pre-specified equivalence margin, and dependent on how unmeasured responses are
treated. The specific measured pairing is not necessary for strong performance
on this benchmark; its biological significance requires further evaluation.
Code, data checksums, results and a supplement are at
\url{https://github.com/ssenge/FlyHash-Connectome}.

\section*{Acknowledgment}

Claude (by Anthropic) was used as a coding and analysis assistant; the author
designed the study.

\footnotesize
\begin{thebibliography}{00}
\setlength{\itemsep}{0pt}

\bibitem{indyk1998} P. Indyk and R. Motwani, ``Approximate nearest neighbors:
towards removing the curse of dimensionality,'' in \emph{Proc. 30th ACM STOC},
1998, pp. 604--613.

\bibitem{charikar2002} M. S. Charikar, ``Similarity estimation techniques from
rounding algorithms,'' in \emph{Proc. 34th ACM STOC}, 2002, pp. 380--388.

\bibitem{dasgupta2017} S. Dasgupta, C. F. Stevens, and S. Navlakha, ``A neural
algorithm for a fundamental computing problem,'' \emph{Science}, vol. 358, no.
6364, pp. 793--796, 2017.

\bibitem{caron2013} S. J. C. Caron, V. Ruta, L. F. Abbott, and R. Axel,
``Random convergence of olfactory inputs in the \emph{Drosophila} mushroom
body,'' \emph{Nature}, vol. 497, pp. 113--117, 2013.

\bibitem{zheng2022} Z. Zheng \textit{et al.}, ``Structured sampling of
olfactory input by the fly mushroom body,'' \emph{Current Biology}, vol. 32,
no. 15, pp. 3334--3349, 2022.

\bibitem{malecns2026} S. Berg, I. R. Beckett, M. Costa, P. Schlegel, M.
Januszewski, E. C. Marin, \textit{et al.}, ``Sexual dimorphism in the complete
\emph{Drosophila} male central nervous system connectome,'' \emph{Cell}, 2026,
doi:10.1016/j.cell.2026.08.015.

\bibitem{strona2014} G. Strona, D. Nappo, F. Boccacci, S. Fattorini, and J.
San-Miguel-Ayanz, ``A fast and unbiased procedure to randomize ecological
matrices with fixed row and column totals,'' \emph{Nature Communications},
vol. 5, art. 4114, 2014.

\bibitem{carstens2015} C. J. Carstens, ``Proof of uniform sampling of binary
matrices with fixed row sums and column sums for the fast Curveball
algorithm,'' \emph{Physical Review E}, vol. 91, art. 042812, 2015.

\bibitem{munch2016} D. M\"unch and C. G. Galizia, ``DoOR 2.0 --- comprehensive
mapping of \emph{Drosophila melanogaster} odorant responses,'' \emph{Scientific
Reports}, vol. 6, art. 21841, 2016.

\bibitem{dasgupta2018} S. Dasgupta, T. C. Sheehan, C. F. Stevens, and S.
Navlakha, ``A neural data structure for novelty detection,'' \emph{Proc. Natl.
Acad. Sci. USA}, vol. 115, no. 51, pp. 13093--13098, 2018.

\end{thebibliography}

\end{document}
